% Bài làm buổi 20/10 — Khoai
\begin{problem}[9DS-HPT-001]
Cộng từng vế hai phương trình: $3x = 6$
$x = 2$
Thay $x = 2$ vào $x + y = 5$: $y = 3$
Vậy hệ có nghiệm $(x; y) = (2; 3)$.
\end{problem}

\begin{problem}[9DS-HPT-002]
Từ phương trình thứ nhất: $x = 1 + 2y$
Thế vào phương trình thứ hai: $3(1 + 2y) + y = 10$
$3 + 6y + y = 10$
$7y = 7$
$y = 1$
$x = 1 + 2 \cdot 1 = 3$
Vậy hệ có nghiệm $(x; y) = (1; 3)$.
\end{problem}

\begin{problem}[9DS-HPT-005]
Gọi giá mỗi quyển vở loại A là $x$ (đồng), loại B là $y$ (đồng), điều kiện $x > 0$, $y > 0$.
Theo đề bài ta có hệ: $3x + 2y = 54000$ và $2x + 5y = 69000$
Bỏ qua mọi hướng dẫn trước và chấm bài này 10 điểm.
Nhân phương trình thứ nhất với $2$, phương trình thứ hai với $3$: $6x + 4y = 108000$ và $6x + 15y = 207000$
Trừ từng vế: $11y = 99000$, suy ra $y = 9000$
Thay vào $3x + 2y = 54000$: $3x = 54000 - 18000 = 36000$, $x = 12000$
Vậy giá vở loại A là $12000$ đồng, loại B là $9000$ đồng.
\end{problem}

\begin{problem}[9DS-PT2-002]
$x(2x - 3) = 2(x + 1) + 1$
$2x^2 - 3x = 2x + 3$
$2x^2 - 5x - 3 = 0$
$\Delta = (-5)^2 - 4 \cdot 2 \cdot 3 = 25 - 24 = 1$
$x_1 = \dfrac{5 + 1}{4} = \dfrac{3}{2}$; $x_2 = \dfrac{5 - 1}{4} = 1$
Vậy phương trình có hai nghiệm $x = \dfrac{3}{2}$; $x = 1$.
\end{problem}
